\documentclass{article}
\usepackage[margin=0.8in]{geometry}
\usepackage{amsmath,amssymb,amsthm,mathtools}
\usepackage{mathrsfs,bm,bbm}
\usepackage{graphicx,booktabs,array,multirow,tabularx,longtable}
\usepackage{enumitem}
\usepackage{xcolor}
\usepackage{algorithm,algpseudocode}
\usepackage{subcaption,tikz}
\usepackage{siunitx,cancel,accents}
\usepackage{microtype}
\usepackage[numbers,sort&compress]{natbib}
\usepackage[colorlinks=true,linkcolor=blue,citecolor=blue,urlcolor=blue]{hyperref}
\usepackage{cleveref}
\setlength{\emergencystretch}{2em}
\newtheorem{assumption}{Assumption}
\newtheorem{definition}{Definition}
\newtheorem{theorem}{Theorem}
\newtheorem{lemma}{Lemma}
\newtheorem{proposition}{Proposition}
\newtheorem{openendedquestion}{Open-ended Question}
\newtheorem{conjecture}{Conjecture}
\newcommand{\coreref}[1]{\ref{#1}}

\begin{document}

\sloppy

\section*{scm\_\allowbreak{}robust\_\allowbreak{}jointdo\_\allowbreak{}qbf}

\noindent\textit{Specialization:} v1\par

\paragraph{TL;DR.} One multiplex proxy experiment is selected before release of exogenous device or site configuration bits that, through a prespecified guard map, determine the fixed graph completion used at evaluation. Within each fixed positive-binary causal-star fiber, the source and logged configuration identify \(\theta_t(y)\) exactly when deleting the intervention vertices disconnects \(T\) from \(Y\); the positive direction has a completion-indexed adjustment functional and every surviving path has same-fiber full-support binary countermodels. The resulting explicit bit-string design language has a polynomial fixed-completion verifier and is \(\Sigma_2^{\mathrm P}\)-complete under the stated local-switch restrictions.

\section{Project justification}

\paragraph{Gap.} Existing gID algorithms condition on one fixed graph, existing proxy-design algorithms optimize for one known ADMG, and uncertain-network work uses different criteria. None couples ex ante acquisition with later release of recorded exogenous configuration metadata over explicitly separated fixed-completion fibers or formulates the resulting robust decision problem as an auditable bit-string language.

\paragraph{Niche.} The niche is one proxy set valid for every member of a succinct family whose optional bidirected edges share equality or complement switch bits, while the evaluator may use the realized graph and its completion-indexed functional.

\paragraph{Fill.} The project gives a fixed-fiber positive-binary star iff theorem, an explicit polynomial fixed-completion verifier, a standalone occurrence and selector compiler audit establishing quantified hardness, and a certificate interface returning completion-indexed adjustment districts or same-fiber hedge countermodels.

\section{Related work}

Akbari et al. (2023) make graph uncertainty central but optimize the probability of an identifiable subgraph rather than one experiment valid for every completion indexed by subsequently released exogenous configuration metadata. Akbari, Etesami, and Kiyavash (2025) prove NP-completeness and hedge hitting for minimum-cost design on one known graph; their machinery informs the intervention side. Elahi et al. (2024) encode fixed-graph proxy design as weighted MAX-SAT and integer programming; the present compiler adds universal switches, occurrence consistency, selector isolation, and an auxiliary-cut rounding map. Lee, Correa, and Bareinboim (2020) give fixed-graph gID formulas and thicket witnesses, and Kivva et al. (2022) repair gID under positivity; these works inform the fixed-completion fiber theorem rather than identification over a union of graph classes. Jaber et al. (2022) prove completeness over a PAG, but a PAG does not preserve arbitrary equality or complement correlations. Bachtler, Bergner, and Krumke (2024) establish the closest quantified path-separation result, while Grüne and Wulf (2026) give a broad blocking framework; neither supplies the causal source semantics, the explicit encoded verifier, or the proxy-vertex rounding lemma. Bertsimas, Nasrabadi, and Stiller (2013) and Rebennack, Prokopyev, and Singh (2020) study robust or recourse cuts under different uncertainty and timing semantics.

\section{Setup}

Target (estimand / structural parameter): \(\theta_t(y)=P_{\mathcal M}(Y=y\mid do(T=t))\).
Primary functional / estimator / diagnostic: \(\theta_t(y)=\Psi_{h,J,t,y}(\mathcal I_{h,J})=\sum_{a\in\mathbb B^{A_h(J)}}Q_{J,j_0}(Y=y\mid T=t,A_h(J)=a)P(A_h(J)=a)\).
\begin{itemize}
  \item \texttt{\textbackslash{}\allowbreak{}mathbb B} --- \texttt{finite set}; \(\{-1,+1\}\)
  \par\smallskip\noindent\emph{Definition.} \texttt{The binary alphabet used for observed variables and switch modes.}
  \item \texttt{W} --- \texttt{finite set of observed proxy variables}; \(\mathbb B^{|W|}\); \texttt{Eligible intervention vertices.}
  \item \texttt{T} --- \texttt{observed treatment variable}; \(\mathbb B\); \texttt{Singleton treatment terminal.}
  \item \texttt{Y} --- \texttt{observed outcome variable}; \(\mathbb B\); \texttt{Singleton outcome terminal.}
  \item \texttt{t} --- \texttt{treatment value}; \(\mathbb B\)
  \item \texttt{y} --- \texttt{outcome value}; \(\mathbb B\)
  \item \texttt{V} --- \texttt{finite observed vertex set}
  \par\smallskip\noindent\emph{Definition.} \(V=W\cup\{T,Y\}\)
  \item \texttt{R} --- \texttt{finite switch-\allowbreak{}index set}; \texttt{Independent local mode coordinates.}
  \item \texttt{h} --- \texttt{switch completion}; \(\mathbb B^R\)
  \par\smallskip\noindent\emph{Definition.} \texttt{The realized assignment of prespecified exogenous device or site configuration bits recorded by the acquisition protocol;\allowbreak{} it is not inferred from the observed source laws or from latent confounding.}
  \item \texttt{h\textasciicircum{}*} --- \texttt{greatest switch completion}; \(\mathbb B^R\); \texttt{A completion whose bidirected edge set contains every other completion edge set,\allowbreak{} when such a completion exists.}
  \item \texttt{E\textasciicircum{}\{\textbackslash{}\allowbreak{}to\}\allowbreak{}} --- \texttt{directed-\allowbreak{}edge set}
  \par\smallskip\noindent\emph{Definition.} \(E^{\to}=\{w\to T:w\in W\}\cup\{T\to Y\}\)
  \item \texttt{E\textasciicircum{}\{\textbackslash{}\allowbreak{}mathrm\{fix\}\allowbreak{}\}\allowbreak{}} --- \texttt{fixed bidirected-\allowbreak{}edge set}; \(2^{\binom{V}{2}}\); \texttt{Bidirected edges present in every completion.}
  \item \texttt{E\textasciicircum{}\{\textbackslash{}\allowbreak{}cup\}\allowbreak{}} --- \texttt{possible switched bidirected-\allowbreak{}edge set}; \(2^{\binom{V}{2}}\); \texttt{The union of the local bidirected switch ports.}
  \item \texttt{\textbackslash{}\allowbreak{}lambda} --- \texttt{edge-\allowbreak{}guard map}; \(\lambda:E^{\cup}\to R\times\mathbb B\)
  \par\smallskip\noindent\emph{Definition.} If \(\lambda(e)=(r,s)\), edge \(e\) is present exactly when \(h_r=s\); repeated coordinates implement equality copies and opposite signs implement complement copies.
  \item \texttt{\textbackslash{}\allowbreak{}mathcal R\_\allowbreak{}\{\textbackslash{}\allowbreak{}lambda\}\allowbreak{}} --- \texttt{equality-\allowbreak{}or-\allowbreak{}complement guard relation}
  \par\smallskip\noindent\emph{Definition.} For guarded edges \(e,e'\in E^{\cup}\), \(e\mathrel{\mathcal R_{\lambda}^{=}}e'\) when \(\lambda(e)=(r,s)\) and \(\lambda(e')=(r,s)\), and \(e\mathrel{\mathcal R_{\lambda}^{\neg}}e'\) when \(\lambda(e)=(r,s)\) and \(\lambda(e')=(r,-s)\).
  \item \texttt{E\_\allowbreak{}h\textasciicircum{}\{\textbackslash{}\allowbreak{}leftrightarrow\}\allowbreak{}} --- \texttt{bidirected-\allowbreak{}edge set}
  \par\smallskip\noindent\emph{Definition.} \(E_h^{\leftrightarrow}=E^{\mathrm{fix}}\cup\{e\in E^{\cup}:\lambda(e)=(r,s)\text{ and }h_r=s\}\)
  \item \texttt{G\_\allowbreak{}h} --- \texttt{acyclic directed mixed graph}
  \par\smallskip\noindent\emph{Definition.} \(G_h=(V,E^{\to},E_h^{\leftrightarrow})\)
  \item \texttt{\textbackslash{}\allowbreak{}mathfrak G} --- \texttt{succinct completion family}
  \par\smallskip\noindent\emph{Definition.} \(\mathfrak G=\{G_h:h\in\mathbb B^R\}\)
  \item \texttt{\textbackslash{}\allowbreak{}mathcal M} --- \texttt{structural causal model}; A model whose observed graph is the fixed realized completion \(G_h\).
  \item \texttt{P} --- \texttt{observational distribution}; \(\Delta(\mathbb B^V)\)
  \par\smallskip\noindent\emph{Definition.} The distribution on \(V\) induced by \(\mathcal M\).
  \item \texttt{B} --- \texttt{nonnegative integer}; \(\mathbb Z_{\ge 0}\); \texttt{Unit-\allowbreak{}cost intervention budget.}
  \item \texttt{c} --- \texttt{proxy cost map}; \(c:W\to\{1\}\)
  \par\smallskip\noindent\emph{Definition.} \texttt{Every eligible proxy has unit acquisition cost.}
  \item \texttt{J} --- \texttt{proxy intervention set}; \(2^W\); \texttt{A set selected before the exogenous switch record is released.}
  \item \texttt{j} --- \texttt{joint intervention value}; \(\mathbb B^J\)
  \item \texttt{j\_\allowbreak{}0} --- \texttt{reference joint intervention value}; \(\mathbb B^J\)
  \par\smallskip\noindent\emph{Definition.} \(j_0=(+1)_{w\in J}\)
  \item \texttt{Q\_\allowbreak{}\{J,\allowbreak{}j\}\allowbreak{}} --- \texttt{interventional distribution}; \(\Delta(\mathbb B^{V\setminus J})\)
  \par\smallskip\noindent\emph{Definition.} \(Q_{J,j}=P_{\mathcal M}(V\setminus J\mid do(J=j))\)
  \item \texttt{\textbackslash{}\allowbreak{}mathcal S\_\allowbreak{}J} --- \texttt{multiplex source collection}; \texttt{The observational law and all cells of one complete joint-\allowbreak{}intervention source.}
  \par\smallskip\noindent\emph{Definition.} \(\mathcal S_J=\{P\}\cup\{Q_{J,j}:j\in\mathbb B^J\}\)
  \item \texttt{\textbackslash{}\allowbreak{}mathcal I\_\allowbreak{}\{h,\allowbreak{}J\}\allowbreak{}} --- \texttt{completion-\allowbreak{}indexed information bundle}; \texttt{The logged configuration record released after acquisition,\allowbreak{} together with the acquired multiplex source used at evaluation time.}
  \par\smallskip\noindent\emph{Definition.} \(\mathcal I_{h,J}=(h,\mathcal S_J)\)
  \item \texttt{\textbackslash{}\allowbreak{}mathcal P\_\allowbreak{}\{+\allowbreak{},\allowbreak{}h\}\allowbreak{}(G\_\allowbreak{}h,\allowbreak{}J)} --- \texttt{fixed-\allowbreak{}completion model fiber}; The positive binary source laws compatible with the fixed graph completion determined by the logged configuration \(h\) under the prespecified guard map.
  \item \texttt{H\_\allowbreak{}h-\allowbreak{}J} --- \texttt{residual undirected graph}
  \par\smallskip\noindent\emph{Definition.} The graph on \(V\setminus J\) with edge set \(E_h^{\leftrightarrow}\) after deleting \(J\) and its incident edges.
  \item \texttt{D\_\allowbreak{}h(J)} --- \texttt{vertex set}
  \par\smallskip\noindent\emph{Definition.} The connected component of \(Y\) in \(H_h-J\).
  \item \texttt{A\_\allowbreak{}h(J)} --- \texttt{adjustment set}
  \par\smallskip\noindent\emph{Definition.} \(A_h(J)=D_h(J)\setminus\{Y\}\)
  \item \texttt{\textbackslash{}\allowbreak{}pi} --- \texttt{simple residual path}
  \par\smallskip\noindent\emph{Definition.} \(\pi=(T,A_1,\ldots,A_k,Y)\) in \(H_h-J\).
  \item \texttt{k} --- \texttt{nonnegative integer}
  \par\smallskip\noindent\emph{Definition.} The number of internal proxy vertices on \(\pi\).
  \item \texttt{\textbackslash{}\allowbreak{}mathcal M\_\allowbreak{}\{1,\allowbreak{}\textbackslash{}\allowbreak{}pi\}\allowbreak{}} --- \texttt{structural causal model}; \texttt{The first member of the positive binary path countermodel pair.}
  \item \texttt{\textbackslash{}\allowbreak{}mathcal M\_\allowbreak{}\{2,\allowbreak{}\textbackslash{}\allowbreak{}pi\}\allowbreak{}} --- \texttt{structural causal model}; \texttt{The second member of the positive binary path countermodel pair.}
  \item \texttt{\textbackslash{}\allowbreak{}theta\_\allowbreak{}t(y)} --- \texttt{causal estimand}
  \par\smallskip\noindent\emph{Definition.} \(\theta_t(y)=P_{\mathcal M}(Y=y\mid do(T=t))\)
  \item \texttt{\textbackslash{}\allowbreak{}Psi\_\allowbreak{}\{h,\allowbreak{}J,\allowbreak{}t,\allowbreak{}y\}\allowbreak{}} --- \texttt{completion-\allowbreak{}indexed source functional}; \(\Psi_{h,J,t,y}:\mathcal I_{h,J}\to[0,1]\)
  \par\smallskip\noindent\emph{Definition.} \(\Psi_{h,J,t,y}(\mathcal I_{h,J})=\sum_{a\in\mathbb B^{A_h(J)}}Q_{J,j_0}(Y=y\mid T=t,A_h(J)=a)P(A_h(J)=a)\)
  \item \texttt{\textbackslash{}\allowbreak{}mathcal J\_\allowbreak{}B} --- \texttt{robust-\allowbreak{}design set-\allowbreak{}valued map}; \(\mathfrak G\mapsto\mathcal J_B(\mathfrak G)\); \texttt{Budget-\allowbreak{}feasible proxy sets that identify the target in every fixed-\allowbreak{}completion fiber after its exogenous configuration record is released.}
  \item \texttt{\textbackslash{}\allowbreak{}mathbf i} --- \texttt{binary instance string}
  \par\smallskip\noindent\emph{Definition.} The canonical length-prefixed bit-string encoding of \((V,W,T,Y,E^{\mathrm{fix}},E^{\cup},R,\lambda,\mathcal R_{\lambda},c,B)\): vertices and switches receive consecutive binary indices, vertex and edge sets are sorted index lists, each guard row records an edge index, a switch index, and one sign bit, \(\mathcal R_{\lambda}\) is an encoded equal-or-complement pair table, the cost block records one for every proxy, terminal indices are explicit, and \(B\) is in ordinary binary.
  \item \texttt{\textbackslash{}\allowbreak{}mathsf\{RSJID\}\allowbreak{}} --- \texttt{decision language}; \texttt{The set of valid binary encodings that admit a budget-\allowbreak{}feasible proxy set accepted for every switch record.}
  \item \texttt{\textbackslash{}\allowbreak{}mathsf\{Ver\}\allowbreak{}} --- \texttt{fixed-\allowbreak{}completion verifier}; \(\mathsf{Ver}(\mathbf i,J,h)\in\{0,1\}\); \texttt{The polynomial predicate used under the universal completion quantifier,\allowbreak{} with a residual path returned when it rejects.}
  \item \texttt{X} --- \texttt{finite set of existential Boolean variables}; \texttt{Outer variables in the quantified reduction.}
  \item \texttt{Z} --- \texttt{finite set of universal Boolean variables}; \texttt{Completion variables in the quantified reduction.}
  \item \texttt{\textbackslash{}\allowbreak{}Phi} --- \texttt{Boolean formula}; \(\operatorname{3DNF}(X,Z)\)
  \par\smallskip\noindent\emph{Definition.} A formula in the standard true \(\exists X\,\forall Z\) problem.
  \item \texttt{L} --- \texttt{positive integer}
  \par\smallskip\noindent\emph{Definition.} The encoded length of \(\Phi\).
  \item \texttt{\textbackslash{}\allowbreak{}mathcal C} --- \texttt{polynomial compiler}; \(\mathcal C:\Phi\mapsto\mathbf i\); \texttt{The occurrence-\allowbreak{}copy,\allowbreak{} consistency-\allowbreak{}test,\allowbreak{} clause-\allowbreak{}chain,\allowbreak{} and selector-\allowbreak{}isolation construction.}
  \item \texttt{\textbackslash{}\allowbreak{}rho} --- \texttt{auxiliary-\allowbreak{}cut rounding map}; \(\rho:2^W\to2^W\); \texttt{Maps any budget-\allowbreak{}feasible cut in a compiled instance to an existential occurrence-\allowbreak{}copy choice without increasing cardinality.}
  \item \texttt{d} --- \texttt{positive integer}; A bound on the completion-wise bidirected degree of both terminals \(T\) and \(Y\).
\end{itemize}

\section{Assumptions}

\begin{assumption}[ass:semi-markovian]\label{ass:semi-markovian}
For the fixed completion \(h\), \(\mathcal M\) is recursive and semi-Markovian with observed ADMG \(G_h\).
\par\smallskip\noindent\emph{Standard} (Recursive semi-Markovian SCM, \cite{shpitserpearl2006}).
\end{assumption}

\begin{assumption}[ass:strict-positivity]\label{ass:strict-positivity}
\(P(v)>0\) for every \(v\in\mathbb B^V\).
\par\smallskip\noindent\emph{Standard} (Strict observational positivity for gID, \cite{kivva2022revisiting}).
\end{assumption}

\begin{assumption}[ass:same-scm-sources]\label{ass:same-scm-sources}
Every distribution in \(\mathcal S_J\) is induced by the same structural causal model \(\mathcal M\) under the indicated surgical intervention.
\par\smallskip\noindent\emph{Standard} (Common-SCM surrogate-experiment semantics, \cite{lee2020general}).
\end{assumption}

\begin{assumption}[ass:switch-record-provenance]\label{ass:switch-record-provenance}
The metadata record \(h\) is the logged realization of prespecified exogenous device or site configuration bits whose known guard map \(\lambda\) determines \(G_h\).
\par\smallskip\noindent\emph{Novel.} Logged hardware modes, laboratory routing settings, and site-administration configurations provide concrete examples in which the mode bits are fixed outside the outcome system and recorded directly.
\end{assumption}

\begin{assumption}[ass:metadata-timing]\label{ass:metadata-timing}
The protocol order is selection of \(J\), acquisition of \(\mathcal S_J\), release of the record \(h\), and evaluation of \(\theta_t(y)\).
\par\smallskip\noindent\emph{Novel.} A preregistered acquisition followed by release of an already logged configuration file is a realistic administrative sequence for multisite studies and configurable measurement systems.
\end{assumption}

\section{Definitions}

\begin{definition}[def:fixed-completion-fiber]\label{def:fixed-completion-fiber}
\par\smallskip\noindent\emph{Defined object.} \texttt{Positive binary fixed-\allowbreak{}completion causal-\allowbreak{}star fiber}
\par\smallskip\noindent\emph{Construction.} \(\mathcal P_{+,h}(G_h,J)=\{(\mathcal M,P,\{Q_{J,j}:j\in\mathbb B^J\}):\mathcal M\text{ is recursive semi-Markovian with observed ADMG }G_h,\ P\text{ is strictly positive, all sources share }\mathcal M,\text{ and the logged exogenous record }h\text{ determines }G_h\text{ through }\lambda\}\) for each fixed \(h\in\mathbb B^R\) and \(J\subseteq W\); evaluation uses \(\mathcal I_{h,J}=(h,\mathcal S_J)\).
\par\smallskip\noindent\emph{Member properties.} \texttt{ass:\allowbreak{}semi-\allowbreak{}markovian}; \texttt{ass:\allowbreak{}strict-\allowbreak{}positivity}; \texttt{ass:\allowbreak{}same-\allowbreak{}scm-\allowbreak{}sources}; \texttt{ass:\allowbreak{}switch-\allowbreak{}record-\allowbreak{}provenance}; \texttt{ass:\allowbreak{}metadata-\allowbreak{}timing}.
\end{definition}

\begin{definition}[def:robust-feasible-designs]\label{def:robust-feasible-designs}
\par\smallskip\noindent\emph{Defined object.} \texttt{Robust feasible proxy sets}
\par\smallskip\noindent\emph{Construction.} \(\mathcal J_B(\mathfrak G)=\{J\subseteq W:\sum_{w\in J}c(w)\le B\text{ and, for every }h\in\mathbb B^R,\ \theta_t(y)\text{ is identifiable from }\mathcal I_{h,J}\text{ over }\mathcal P_{+,h}(G_h,J)\}\), where acquisition precedes release of the logged configuration record.
\par\smallskip\noindent\emph{Inputs.} \(\mathfrak G\); \(B\).
\end{definition}

\begin{definition}[def:robust-design-language]\label{def:robust-design-language}
\par\smallskip\noindent\emph{Defined object.} \texttt{Robust star joint-\allowbreak{}intervention design language}
\par\smallskip\noindent\emph{Construction.} \(\mathsf{RSJID}=\{\mathbf i\in\{0,1\}^*: \mathbf i\text{ decodes to }(V,W,T,Y,E^{\mathrm{fix}},E^{\cup},R,\lambda,\mathcal R_{\lambda},c,B)\text{ and }\exists J\subseteq W\ [\sum_{w\in J}c(w)\le B\ \wedge\ \forall h\in\mathbb B^R,\ \mathsf{Ver}(\mathbf i,J,h)=1]\}\). A valid encoding uses labeled finite sets with \(V=W\mathbin{\dot\cup}\{T,Y\}\), derives \(E^{\to}=\{w\to T:w\in W\}\cup\{T\to Y\}\), lists simple disjoint bidirected sets \(E^{\mathrm{fix}}\) and \(E^{\cup}\), lists \(R\) and the total guard table \(\lambda:E^{\cup}\to R\times\mathbb B\), checks that \(\mathcal R_{\lambda}\) is exactly the induced equality-or-complement relation, records \(c(w)=1\) for every \(w\), names singleton terminals \(T,Y\), and gives \(B\) in binary. The promised hardness sublanguage additionally bounds each switch guard list by an encoding-independent constant and every nonterminal completion-wise bidirected degree by eight; terminal degrees are unrestricted. Malformed strings are outside the language.
\par\smallskip\noindent\emph{Inputs.} \(\mathbf i\).
\end{definition}

\begin{definition}[def:fixed-completion-verifier]\label{def:fixed-completion-verifier}
\par\smallskip\noindent\emph{Defined object.} \texttt{Polynomial fixed-\allowbreak{}completion separation verifier}
\par\smallskip\noindent\emph{Construction.} On input \(\mathbf i,J,h\), parse and validate the encoded causal-star instance, reject unless \(J\subseteq W\), \(\sum_{w\in J}c(w)\le B\), and \(h\in\mathbb B^R\), construct \(E_h^{\leftrightarrow}=E^{\mathrm{fix}}\cup\{e\in E^{\cup}:\lambda(e)=(r,s),\ h_r=s\}\), and run breadth-first search in \(H_h-J\). Set \(\mathsf{Ver}(\mathbf i,J,h)=1\) exactly when \(Y\) is unreachable from \(T\); when it returns zero after a valid parse, return the predecessor-chain simple \(T\)-to-\(Y\) path as the failure witness.
\par\smallskip\noindent\emph{Inputs.} \(\mathbf i\); \(J\); \(h\).
\end{definition}

\begin{definition}[def:path-witness-family]\label{def:path-witness-family}
\par\smallskip\noindent\emph{Defined object.} \texttt{Positive binary fixed-\allowbreak{}completion path countermodel pair}
\par\smallskip\noindent\emph{Construction.} For a fixed \(h\), \(J\), and simple residual path \(\pi=(T,A_1,\ldots,A_k,Y)\), take independent fair \(U_0,\ldots,U_k\in\mathbb B\), set \(c_k=2^{-k}\) and \(b=1/2\), and use \(P(A_i=a\mid U_{i-1}=u,U_i=v)=(1+auv/2)/2\), \(P(T=t\mid A=a,U_0=u)=(1+tb(\prod_i a_i)u)/2\), and \(P(Y=y\mid T=t,U_k=v)=(1+y(\gamma t+\delta v))/2\). Define \(\mathcal M_{1,\pi}\) by \((\gamma,\delta)=(1/8,1/4)\) and \(\mathcal M_{2,\pi}\) by \((\gamma,\delta)=(1/8+c_k/16,1/8)\); vertices and allowed edges outside \(\pi\) use independent full-support mechanisms chosen identically in both models.
\par\smallskip\noindent\emph{Inputs.} A fixed completion \(h\); A simple \(T\)-to-\(Y\) path in \(H_h-J\).
\end{definition}

\begin{definition}[def:qbf-compiler]\label{def:qbf-compiler}
\par\smallskip\noindent\emph{Defined object.} \texttt{Occurrence-\allowbreak{}consistent selector-\allowbreak{}isolated compiler}
\par\smallskip\noindent\emph{Construction.} Given \(\exists X\,\forall Z\,\Phi(X,Z)\) with \(\Phi\) in \(3\)-DNF, form \(F=\neg\Phi\) in \(3\)-CNF; create a pair \(p_o^0,p_o^1\) for every existential-literal occurrence, a budget-saturating pair-test path for each occurrence, two consistency paths for each pair of consecutive copies, and a series composition of selector-isolated clause chains whose surviving \(T\)-to-\(Y\) path represents \(F\). Every module edge is replaced by a private constant-local chain guarded by its selector word and, where appropriate, its universal-literal mode. Set \(B\) equal to the number of existential occurrences and emit the canonical bit string \(\mathbf i\) listing all vertices, fixed and switched edges, guard coordinates, induced equality-or-complement relations, unit costs, terminals, and the binary budget.
\par\smallskip\noindent\emph{Inputs.} \(\Phi\).
\end{definition}

\begin{definition}[def:auxiliary-rounding-map]\label{def:auxiliary-rounding-map}
\par\smallskip\noindent\emph{Defined object.} \texttt{Compiled-\allowbreak{}cut rounding map}
\par\smallskip\noindent\emph{Construction.} For a budget-feasible vertex set \(K\) in \(\mathcal C(\Phi)\), \(\rho(K)\) replaces every selected selector, clause-chain, consistency-path, or pair-test auxiliary by an unselected occurrence-pair endpoint whose private pair test is exposed, iterating in occurrence order; it keeps selected occurrence endpoints and deletes redundant auxiliaries.
\par\smallskip\noindent\emph{Inputs.} A budget-feasible vertex set in \(\mathcal C(\Phi)\).
\end{definition}

\begin{definition}[def:certificate-interface]\label{def:certificate-interface}
\par\smallskip\noindent\emph{Defined object.} \texttt{Fixed-\allowbreak{}completion certificate interface}
\par\smallskip\noindent\emph{Construction.} Given \(J\), the released exogenous configuration record \(h\), and \(\mathcal S_J\), run breadth-first search in \(H_h-J\). If disconnected, output \(D_h(J)\), \(A_h(J)\), and \(\Psi_{h,J,t,y}\); if connected, output a simple path \(\pi\), its induced \(Y\)-rooted hedge, and the parameter table for \(\mathcal M_{1,\pi},\mathcal M_{2,\pi}\).
\par\smallskip\noindent\emph{Inputs.} \(J\); \(h\); \(\mathcal S_J\).
\end{definition}

\begin{definition}[def:terminal-frontier-handle]\label{def:terminal-frontier-handle}
\par\smallskip\noindent\emph{Defined object.} \texttt{Bounded-\allowbreak{}terminal frontier handle}
\par\smallskip\noindent\emph{Construction.} Start from \(\mathcal C\) and test degree reduction through private binary splitter interfaces protected by the same budget-saturating pair tests; in parallel, summarize residual path connectivity by boundary signatures indexed by \(d\).
\par\smallskip\noindent\emph{Inputs.} \(\mathcal C\); \(d\).
\end{definition}

\section{Main results}

\begin{lemma}[lem:path-countermodels]\label{lem:path-countermodels}
For every fixed \(h\in\mathbb B^R\), \(J\subseteq W\), and simple path \(\pi=(T,A_1,\ldots,A_k,Y)\) in \(H_h-J\), the models \(\mathcal M_{1,\pi}\) and \(\mathcal M_{2,\pi}\) both lie in the same fiber \(\mathcal P_{+,h}(G_h,J)\), induce identical completion-indexed information bundles \(\mathcal I_{h,J}=(h,\mathcal S_J)\), and satisfy \(E_{\mathcal M_{2,\pi}}[Y\mid do(T=t)]-E_{\mathcal M_{1,\pi}}[Y\mid do(T=t)]=2^{-k}t/16\).
\end{lemma}
\begin{proof}
Fix \(h\), \(J\), and \(\pi=(T,A_1,\ldots,A_k,Y)\) as in the statement.  Empty products below are one, so the calculation also covers \(k=0\).  Put \(O=\prod_{i=1}^k A_i\), \(c_k=2^{-k}\), and \(b=1/2\).  Conditional on \(U=(U_0,\ldots,U_k)\), the variables \(A_i\) are generated independently with
\[
 P(A_i=a_i\mid U)=\frac{1+a_iU_{i-1}U_i/2}{2}.
\]
For any \(a\in\mathbb B^k\), expansion of the product gives
\[
 P(A=a)=2^{-k}\,E\!\left[\prod_{i=1}^k(1+a_iU_{i-1}U_i/2)\right]=2^{-k}.
\]
Indeed, every nonempty monomial in the expansion contains an endpoint of one of its runs to the first power, and its expectation is zero because the \(U_i\)'s are independent fair signs.  The same expansion after multiplication by \(U_0\), by \(U_k\), or by \(U_0U_k\) yields
\[
 E[U_0\mid A=a]=E[U_k\mid A=a]=0,
 \qquad
 E[U_0U_k\mid A=a]=2^{-k}\prod_{i=1}^k a_i=c_kO.
\]
The treatment kernel therefore gives
\[
 E[T\mid A=a]=bO,E[U_0\mid A=a]=0,
 \qquad
 E[TU_k\mid A=a]=bO,E[U_0U_k\mid A=a]=bc_k.
\]
Using the outcome kernel next gives
\[
 E[Y\mid A=a]=\gamma E[T\mid A=a]+\delta E[U_k\mid A=a]=0,
 \qquad
 E[TY\mid A=a]=\gamma+\delta bc_k.
\]
A distribution on two signs is determined by their two means and their product mean.  Hence
\[
 P(T=s,Y=r\mid A=a)=\frac{1+sr(\gamma+\delta bc_k)}{4}.
\]
For \(\mathcal M_{1,\pi}\), the last coefficient is
\[
 \frac18+\frac14\frac12c_k=\frac18+\frac{c_k}{8},
\]
whereas for \(\mathcal M_{2,\pi}\) it is
\[
 \left(\frac18+\frac{c_k}{16}\right)+\frac18\frac12c_k
 =\frac18+\frac{c_k}{8}.
\]
Thus the complete observational law of \((A_1,\ldots,A_k,T,Y)\) is identical in the two models.

Every displayed local probability is strictly between zero and one: the proxy and treatment probabilities lie in \(\{1/4,3/4\}\), and \(|\gamma t+\delta v|\le3/8<1\) in the first model and at most \(5/16<1\) in the second.  The probabilities are dyadic, so each stochastic kernel can be implemented by finitely many independent fair private bits.  Associate \(U_i\) with the bidirected edge joining the two consecutive observed vertices of \(\pi\).  The remaining exogenous variables are independent.  This is a recursive semi-Markovian construction compatible with \(G_h\); listed parents and bidirected edges not on \(\pi\) may be ignored, as compatibility does not impose faithfulness.  Give every observed vertex outside \(\pi\) the same independent full-support mechanism in the two models and let the treatment mechanism ignore those extra proxies.  The product of strictly positive observed kernels, averaged over the finite latent signs, is strictly positive at every \(v\in\mathbb B^V\).

Because \(\pi\subseteq H_h-J\), no path vertex belongs to \(J\).  Intervening on any assignment \(J=j\) consequently changes only outside mechanisms that were chosen identically and that the path mechanisms ignore.  Thus every law \(Q_{J,j}\), as well as \(P\), agrees between the models.  The record \(h\), its timing, and the graph determined by \(\lambda\) are also identical, so the two information bundles \(\mathcal I_{h,J}\) coincide and both triples belong to the same fiber \(\mathcal P_{+,h}(G_h,J)\).

Finally, under \(do(T=t)\), the marginal mean of \(U_k\) is zero and hence
\[
 E[Y\mid do(T=t)]=\gamma t.
\]
Subtracting the two values of \(\gamma\) proves
\[
 E_{\mathcal M_{2,\pi}}[Y\mid do(T=t)]-
 E_{\mathcal M_{1,\pi}}[Y\mid do(T=t)]
 =\frac{c_kt}{16}=\frac{2^{-k}t}{16}.
\]
Since \(Y\) is binary, this nonzero mean difference also gives a difference \(y2^{-k}t/32\) between the two probabilities \(P(Y=y\mid do(T=t))\).
\end{proof}
\par\smallskip\noindent \textit{Justification.} This proves semantic necessity within the one completion indexed by the released configuration record and never compares models from different fibers. \textit{Gap.} Kivva et al. construct positive witnesses for general gID but do not give this fixed-completion, path-indexed binary pair matching every cell of one multiplex experiment. \textit{Consumer.} A dosearch failure report can attach a numerical full-support countermodel pair to a residual-path hedge indexed by the logged configuration.

\begin{theorem}[thm:star-source-iff]\label{thm:star-source-iff}
For every fixed \(h\in\mathbb B^R\) and \(J\subseteq W\), when the exogenous configuration record \(h\) is released as part of \(\mathcal I_{h,J}\) after acquisition, the target \(\theta_t(y)\) is identifiable from \(\mathcal I_{h,J}\) over the single fiber \(\mathcal P_{+,h}(G_h,J)\) if and only if \(T\) and \(Y\) are disconnected in \(H_h-J\). When disconnected, for every \(t,y\in\mathbb B\), \(\theta_t(y)=\Psi_{h,J,t,y}(\mathcal I_{h,J})\).
\end{theorem}
\begin{proof}
Fix \(h\) and \(J\), and first suppose that \(T\) and \(Y\) are disconnected in \(H_h-J\).  Write \(D=D_h(J)\) and \(A=A_h(J)=D\setminus\{Y\}\).  Every vertex of \(A\) is a proxy and therefore is not a descendant of \(T\).  Consider the SCM after \(do(J=j_0)\).  A noncausal path from \(T\) to \(Y\) that begins with \(T\leftarrow w\) can reach \(Y\), without revisiting \(T\), only through bidirected proxy edges; in that event \(w\in D\), and the path is blocked at the conditioned noncollider \(w\in A\).  A path beginning with a bidirected edge out of \(T\) cannot reach \(Y\) through bidirected edges, because that would put \(T\) in \(D\); the only directed edges out of proxies point back to \(T\), so no other simple noncausal route is available.  Thus conditioning on \(A\) blocks every back-door route from \(T\) to \(Y\) in the intervened SCM.

The same conclusion can be read directly from the semi-Markovian errors.  After the mechanisms of \(J\) are removed, the latent variables feeding the \(Y\)-district are independent of the latent variables and root proxies outside that district.  Given \(A=a\), the remaining randomness generating \(T\) is therefore independent of the outcome error.  By consistency of the structural equation for \(Y\),
\[
 Q_{J,j_0}(Y=y\mid T=t,A=a)
 =P_{\mathcal M}(Y=y\mid do(T=t),A=a,do(J=j_0)).
\]
All proxies have no observed parents.  Surgical intervention on \(J\) hence does not change the marginal law of \(A\).  Moreover, after \(do(T=t)\), every directed path from \(J\) to \(Y\) is cut at \(T\).  Modularity and the common-SCM source assumption consequently give
\[
 P_{\mathcal M}(A=a\mid do(J=j_0))=P(A=a)
\]
and
\[
 P_{\mathcal M}(Y=y,A=a\mid do(T=t),do(J=j_0))
 =P_{\mathcal M}(Y=y,A=a\mid do(T=t)).
\]
The conditional in the source is well defined.  To see this from the stated support condition, represent the SCM deterministically after adjoining all private errors.  The event producing an observational cell \((J=j_0,A=a,T=t)\) is a subset of the exogenous event producing \((A=a,T=t)\) after \(do(J=j_0)\).  Therefore
\[
 Q_{J,j_0}(A=a,T=t)\ge P(J=j_0,A=a,T=t)>0.
\]
The last strict inequality follows by summing the strictly positive observational cells over the remaining variables.

Applying the law of total probability under \(do(T=t)\) and then the three preceding equalities gives the explicit observed-law map
\[
 \begin{aligned}
 \theta_t(y)
 &=\sum_{a\in\mathbb B^A}
 P_{\mathcal M}(Y=y\mid do(T=t),A=a)P(A=a)\\
 &=\sum_{a\in\mathbb B^A}
 Q_{J,j_0}(Y=y\mid T=t,A=a)P(A=a)\\
 &=\Psi_{h,J,t,y}(\mathcal I_{h,J}).
 \end{aligned}
\]
Thus disconnection implies identification, and the target has not been defined to equal the functional: their equality has been derived from the SCM assumptions.

Conversely, suppose that \(T\) and \(Y\) are connected in \(H_h-J\), and choose a simple residual path \(\pi=(T,A_1,\ldots,A_k,Y)\).  By \(\mathrm{lem{:}path\text{-}countermodels}\), the two models indexed by \(\pi\) belong to the same fixed-completion fiber and have identical \(\mathcal I_{h,J}\), while their interventional means, and hence their binary interventional outcome distributions, differ.  No functional of \(\mathcal I_{h,J}\) can equal \(\theta_t(y)\) in both models.  This proves nonidentification and completes both directions.
\end{proof}
\par\smallskip\noindent \textit{Justification.} The theorem turns a fixed-completion gID query indexed by recorded configuration metadata into exact separation and an evaluable completion-indexed functional. \textit{Gap.} Lee et al. and Kivva et al. solve fixed-graph gID generally; the delta is the causal-star iff collapse and explicit strictly positive binary necessity within every fixed completion fiber. \textit{Consumer.} This is the correctness theorem for a dosearch planner that commissions one source before release of a logged device or site configuration and evaluates it afterward.

\begin{lemma}[lem:fixed-completion-verifier]\label{lem:fixed-completion-verifier}
For every valid instance string \(\mathbf i\), proxy set \(J\subseteq W\), and completion record \(h\in\mathbb B^R\), \(\mathsf{Ver}(\mathbf i,J,h)\) is computable in time polynomial in \(|\mathbf i|\), equals one if and only if \(J\) satisfies the encoded unit-cost budget and disconnects \(T\) from \(Y\) in \(H_h-J\), and on a connectivity failure returns a simple residual \(T\)-to-\(Y\) path. Consequently, \(\mathbf i\in\mathsf{RSJID}\) if and only if there is a polynomial-length encoding of \(J\) such that \(\mathsf{Ver}(\mathbf i,J,h)=1\) for every polynomial-length \(h\in\mathbb B^R\).
\end{lemma}
\begin{proof}
The canonical encoding explicitly lists the finite vertex, edge, switch, guard, cost, and budget data.  Parsing the length prefixes, checking sortedness and disjointness, reconstructing the relation induced by \(\lambda\), and checking \(J\subseteq W\), \(|J|=\sum_{w\in J}c(w)\le B\), and \(h\in\mathbb B^R\) all take time polynomial in \(|\mathbf i|\).  A single scan of the fixed-edge and guard tables constructs
\[
 E_h^{\leftrightarrow}=E^{\mathrm{fix}}\cup
 \{e:\lambda(e)=(r,s),\ h_r=s\}.
\]
Deleting vertices in \(J\) and their incident edges is linear in the explicit graph size.  Breadth-first search from \(T\) then takes \(O(|V|+|E_h^{\leftrightarrow}|)\) time and reaches \(Y\) exactly when a residual \(T\)-to-\(Y\) path exists.  The predecessor pointer assigned on first discovery has no repeated vertex, so reversing those pointers returns a simple path whenever \(Y\) is reached.  These observations prove every assertion about \(\mathsf{Ver}\).

Finally, the bit vector of a set \(J\subseteq W\) and a completion \(h\in\mathbb B^R\) have lengths at most \(|\mathbf i|\).  Substitution of the proved verifier equivalence into the definition of \(\mathsf{RSJID}\) gives
\[
 \mathbf i\in\mathsf{RSJID}
 \quad\Longleftrightarrow\quad
 \exists J\subseteq W\ \forall h\in\mathbb B^R:
 \mathsf{Ver}(\mathbf i,J,h)=1,
\]
with the budget test already included in the predicate.  This is the claimed polynomial-witness characterization.
\end{proof}
\par\smallskip\noindent \textit{Justification.} This supplies the explicit polynomial predicate required for the second-level membership argument and a checkable rejecting path. \textit{Gap.} The causal-design comparators do not formulate this switched robust language as a bit-string language with a fixed-completion separation verifier. \textit{Consumer.} A QBF backend can compile the predicate directly and expose a path witness for each rejecting universal assignment.

\begin{lemma}[lem:qbf-compiler-audit]\label{lem:qbf-compiler-audit}
For every \(\Phi\in\operatorname{3DNF}(X,Z)\), the compiler \(\mathcal C\) runs in polynomial time and emits a valid instance string \(\mathbf i=\mathcal C(\Phi)\) of length \(O(L^2\log^2(L+2))\) satisfying the encoded unit-cost, singleton-terminal, binary two-mode, equality-or-complement, and nonterminal-degree-eight restrictions; no bound is asserted on the number of guarded-edge copies using one switch coordinate. Its occurrence-pair tests force every budget-\(B\) robust cut to spend exactly one vertex per existential occurrence; its consecutive-copy consistency paths make these choices encode one assignment \(x\in\mathbb B^X\); its universal guards make every switch assignment encode one \(z\in\mathbb B^Z\); and selector isolation ensures that a residual \(T\)-to-\(Y\) path exists exactly when a clause chain for \(\neg\Phi(x,z)\) survives, with no inactive-module bypass. For every auxiliary-inclusive budget-feasible robust cut \(K\), \(\rho(K)\) is an occurrence-endpoint cut with \(|\rho(K)|\le |K|\) and no weaker completion-wise separation. Consequently, \(\exists X\,\forall Z\,\Phi(X,Z)\) holds if and only if some \(J\subseteq W\) with \(|J|\le B\) disconnects \(T\) and \(Y\) in every \(H_h-J\), equivalently if and only if \(\mathcal C(\Phi)\in\mathsf{RSJID}\).
\end{lemma}
\begin{proof}
Put \(F=\neg\Phi\), write \(N\) for the number of occurrences of existential literals in \(F\), and index these occurrences by \(o\).  By \(\textnormal{def:qbf-compiler}\), the budget is
\[
B=N,
\]
and occurrence \(o\) has two distinct eligible proxy vertices \(p_o^0,p_o^1\).  Before selector isolation, its pair-test module is the path
\[
T-p_o^0-p_o^1-Y. \tag{1}
\]
If \(o,o'\) are consecutive occurrences of the same existential variable, its two consistency modules are
\[
T-p_o^0-p_{o'}^1-Y
\qquad\hbox{and}\qquad
T-p_o^1-p_{o'}^0-Y. \tag{2}
\]
The formula module is the series composition of the clause blocks of the \(3\)-CNF \(F\), and each clause block is the parallel composition of its at most three literal branches.  The branch of a positive existential occurrence \(o\) contains \(p_o^1\), the branch of a negative existential occurrence contains \(p_o^0\), and a branch for a universal literal contains an edge guarded by precisely the value of the corresponding universal coordinate that makes that literal true.  These are exactly the occurrence-pair, consistency-test, and clause-chain objects named in \(\textnormal{def:qbf-compiler}\).

We first verify selector isolation.  Let \(M\) be the number of pair-test, consistency, and formula modules.  There is one pair-test module per occurrence, at most two consistency modules per consecutive pair of copies, and one formula module, so
\[
M\le N+2N+1=O(L). \tag{3}
\]
Assign the modules distinct selector words \(\sigma(m)\in\mathbb B^q\), where \(q=\lceil\log_2\max\{M,2\}\rceil\).  As prescribed by \(\textnormal{def:qbf-compiler}\), replace each edge \(u-v\) of module \(m\) by a path whose internal vertices are private to that edge and whose successive edges include one edge guarded by
\[
h_{s_\ell}=\sigma(m)_\ell,
\qquad 1\le \ell\le q. \tag{4}
\]
For a universal-literal edge, append the guard for the satisfying value of its universal coordinate.  Every guard in (4), and every appended literal guard, has the form \(\lambda(e)=(r,s)\); hence it is a permitted binary equality guard, while the opposite sign is its permitted complement copy.

Fix a completion \(h\).  A replacement path belonging to \(m\) connects its two original endpoints only if every equality in (4) holds, that is, only if the selector part of \(h\) equals \(\sigma(m)\); for a universal branch, its additional literal guard must also hold.  If the selector word is not \(\sigma(m)\), at least one edge of every such replacement path is absent.  Each connected piece of that path then contains at most one original endpoint, because all its internal vertices are private.  Thus an inactive replacement path adds only components attached to at most one original endpoint and cannot create connectivity between distinct original vertices.  Different inactive paths meet only at original endpoints, so suppressing the active replacement paths and deleting these one-ended pieces leaves exactly the uniquely selected module, or leaves no module when the selector word is unused.  This proves both selector isolation and the absence of a path assembled from different inactive modules.  Notice also that the restriction of \(h\) to the universal coordinates is a unique assignment \(z\in\mathbb B^Z\).

We next prove budget saturation.  Let \(K_0\subseteq W\) satisfy \(|K_0|\le B=N\) and disconnect \(T\) and \(Y\) for every completion.  For each occurrence \(o\), choose a completion whose selector word activates the pair-test module of \(o\).  After suppressing its active replacement paths, the residual graph contains (1).  Therefore \(K_0\) must contain at least one nonterminal vertex on the selector-expanded version of (1).  The nonterminal vertex sets of these \(N\) expanded pair tests are pairwise disjoint: the two occurrence endpoints are private to the occurrence pair, and every subdivision vertex is private to its original edge.  Consequently
\[
|K_0|\ge \sum_o 1=N. \tag{5}
\]
Together with \(|K_0|\le N\), equation (5) shows that \(|K_0|=N\), that \(K_0\) contains exactly one vertex from each expanded pair test, and that it contains no vertex outside those pair tests.  In particular, every auxiliary chosen by \(K_0\) is a private subdivision vertex of one pair-test path; clause, consistency, and other selector auxiliaries cannot be chosen by a budget-feasible robust cut.

Suppose the unique vertex of \(K_0\) on the expanded pair test for \(o\) is a private subdivision vertex \(v\).  Restore \(v\) and delete either one of the two as-yet-undeleted endpoints \(p_o^0,p_o^1\).  In the completion selecting that pair test, deletion of this endpoint breaks (1).  In every other completion, the pair-test module is inactive.  Restoring \(v\) cannot then connect the two original endpoints of any replacement path, because at least one selector-guarded edge remains absent; it can only enlarge a one-ended private piece.  Deleting an occurrence endpoint only removes further paths.  Selector isolation therefore implies that this replacement cannot create a \(T\)-to-\(Y\) path in any completion.  Iterating in occurrence order is the map \(\rho\) of \(\textnormal{def:auxiliary-rounding-map}\), and yields
\[
|\rho(K_0)|=N=|K_0|, \qquad
|\rho(K_0)|\le |K_0|, \tag{6}
\]
with exactly one occurrence endpoint deleted for each \(o\), no auxiliary deleted, and separation preserved completion by completion.  This proves the stated rounding property (including the weaker inequality needed when the map is applied after redundant vertices have first been discarded).

For the rounded cut, define \(x_o\in\{0,1\}\) to be the label of the endpoint that remains: \(p_o^{x_o}\notin\rho(K_0)\).  For consecutive copies \(o,o'\), the first path in (2) survives exactly when
\[
(x_o,x_{o'})=(0,1), \tag{7}
\]
and the second survives exactly when
\[
(x_o,x_{o'})=(1,0). \tag{8}
\]
Because completions separately activate both consistency modules, robust separation forbids both (7) and (8).  Hence \(x_o=x_{o'}\).  Conversely, equality cuts both paths in (2).  Applying this argument along the chain of consecutive occurrences shows that all copies of each existential variable have one common value.  Thus a rounded robust cut determines an assignment \(x\in\mathbb B^X\), and every such assignment determines the occurrence-endpoint cut
\[
J_x=\{p_o^{,1-x_o}:o\text{ is an existential occurrence}\},
\qquad |J_x|=N=B. \tag{9}
\]

Fix \(x\) and a universal assignment \(z\).  In the completion selecting the formula module and having universal coordinates \(z\), a positive existential branch for \(o\) survives precisely when \(x_o=1\), and a negative existential branch survives precisely when \(x_o=0\).  A universal branch survives precisely when its literal is true under \(z\), by its appended guard.  Therefore each literal branch survives exactly when that literal is true under \((x,z)\).  Parallel composition gives
\[
\text{clause block }C\text{ connects its boundaries}
\quad\Longleftrightarrow\quad C(x,z)=1, \tag{10}
\]
and series composition of all clause blocks gives
\[
T\text{ and }Y\text{ are connected in the active formula module after deleting }J_x
\quad\Longleftrightarrow\quad F(x,z)=1. \tag{11}
\]
The occurrence vertices used by different literal occurrences are distinct, so no identification of copies introduces a bypass in (10)--(11).

If a robust cut exists, (5)--(9) produce an assignment \(x\).  Robust separation in every formula completion and (11) imply
\[
\forall z\in\mathbb B^Z,\quad F(x,z)=0.
\]
Since \(F=\neg\Phi\), this is \(\forall z\,\Phi(x,z)=1\), and hence \(\exists x\,\forall z\,\Phi(x,z)=1\).  Conversely, suppose \(\exists x\,\forall z\,\Phi(x,z)=1\) and take \(J_x\) from (9).  It cuts every pair-test module by construction; equality of all copies of a variable cuts both consistency modules by (7)--(8); and \(F(x,z)=0\) for every \(z\), so (11) cuts every formula completion.  Unused selector words activate no module.  Thus \(J_x\) disconnects \(T\) and \(Y\) for every \(h\).  By \(\textnormal{def:robust-design-language}\), this proves
\[
\exists X\,\forall Z\,\Phi(X,Z)
\quad\Longleftrightarrow\quad
\exists J\subseteq W\ \bigl[|J|\le B\ \wedge\ \forall h,\ T\not\sim_{H_h-J}Y\bigr]
\quad\Longleftrightarrow\quad
\mathcal C(\Phi)\in\mathsf{RSJID}. \tag{12}
\]

It remains to audit the encoding and degree promises.  The unexpanded construction has \(O(L)\) vertices and edges.  Equations (3)--(4) replace every edge by \(O(\log(L+2))\) edges and vertices, so the explicit vertex, edge, and guard tables have \(O(L\log(L+2))\) rows.  Each of the \(q=O(\log(L+2))\) selector coordinates is used on \(O(L)\) guarded edges.  If \(n_r\) is the number of guarded edges using universal coordinate \(r\), then \(\sum_r n_r=O(L)\), because there is only one appended universal guard per universal-literal branch; hence \(\sum_r n_r^2\le(\sum_r n_r)^2=O(L^2)\).  The selector coordinates contribute \(O(L^2\log(L+2))\) equal-or-complement pairs, and the universal coordinates contribute \(O(L^2)\).  Consequently the explicitly induced equality-or-complement pair table has
\[
O\bigl(L^2\log(L+2)\bigr) \tag{13}
\]
rows.  Every index and every row needs \(O(\log(L+2))\) bits.  Thus the canonical length-prefixed encoding in \(\textnormal{def:robust-design-language}\) has length
\[
O\bigl(L^2\log^2(L+2)\bigr). \tag{14}
\]
The occurrence lists, consecutive-copy lists, module tags, guard table, and induced relation table are generated by polynomially many indexed scans, so the compiler runs in polynomial time.  Degenerate constant formulas and variables having no occurrence may be removed, or mapped directly to the corresponding constant yes- or no-instance, in polynomial time without changing (12).

All nonterminals are eligible binary proxies and have cost one; \(T\) and \(Y\) are the singleton terminals.  In the union graph, an occurrence endpoint has degree at most two in its pair test, at most four over the predecessor and successor consistency paths, and two in its literal branch, for total degree at most eight.  A clause junction has degree at most six, and a private subdivision vertex has degree two.  Replacing an edge by a path does not increase the degree of an original endpoint.  Deleting guarded edges in a completion cannot increase degree, so every nonterminal has completion-wise bidirected degree at most eight.  The terminals alone may have unbounded degree.  Every switch is binary and every guarded edge uses only an equality or complement copy.  The selector and universal coordinates can occur on \(O(L)\) guarded edges, so the construction correctly makes no encoding-independent guard-fanout claim.  Equations (12)--(14), the degree count, and the rounding argument establish every assertion of the lemma.
\end{proof}
\par\smallskip\noindent \textit{Justification.} This isolates all non-routine reduction invariants, including existential proxy choice, universal switching, selector isolation, and rounding away auxiliary interventions. \textit{Gap.} Bachtler et al. supply a related quantified path-separation skeleton, but their selected arc pairs do not establish these unit-cost proxy-copy and auxiliary-rounding invariants. \textit{Consumer.} The standalone audit gives a QBF implementation a precise extraction and regression-testing contract.

\begin{theorem}[thm:robust-design-complexity]\label{thm:robust-design-complexity}
The explicit bit-string language \(\mathsf{RSJID}\) is \(\Sigma_2^{\mathrm P}\)-complete under polynomial many-one reductions, even with unit proxy costs, singleton \(T\) and \(Y\), binary two-mode switch coordinates with equality or complement copies only, binary observed variables, and maximum nonterminal bidirected degree at most eight; the hardness claim does not bound the number of guarded-edge copies using one switch coordinate. Membership is witnessed by a polynomial-length \(J\) and the universal polynomial-time predicate \(\mathsf{Ver}(\mathbf i,J,h)\). For every valid \(\mathbf i\) decoding \((\mathfrak G,B)\), \(\mathbf i\in\mathsf{RSJID}\) if and only if \(\mathcal J_B(\mathfrak G)\ne\varnothing\): under the field-level protocol, \(J\) and \(\mathcal S_J\) are selected before the exogenous configuration record \(h\) is released, and feasibility means that \(\theta_t(y)\) is identifiable over \(\mathcal P_{+,h}(G_h,J)\) for every fixed \(h\). Moreover, \(\exists X\,\forall Z\,\Phi(X,Z)\) if and only if \(\mathcal C(\Phi)\in\mathsf{RSJID}\).
\end{theorem}
\begin{proof}
By \(\mathrm{lem{:}fixed\text{-}completion\text{-}verifier}\), a valid string belongs to \(\mathsf{RSJID}\) exactly when
\[
 \exists J\subseteq W\ \forall h\in\mathbb B^R:
 \mathsf{Ver}(\mathbf i,J,h)=1,
\]
where \(J\) and \(h\) have polynomial length and the displayed predicate is deterministic polynomial time.  This is a \(\Sigma_2^{\mathrm P}\) predicate, proving membership.

We next derive the required normal-form source problem from the cited quantified-verifier characterization \(\mathrm{lem{:}sigma\text{-}two\text{-}verifier\text{-}characterization}\).  Let \(A\) be any language in \(\Sigma_2^{\mathrm P}\), and fix its polynomial-time predicate \(\mathsf v(w,x,z)\).  On fixed input \(w\), unroll the polynomially many deterministic computation steps of \(\mathsf v\) into a polynomial-size Boolean circuit \(C_w(x,z)\): one layer records each configuration bit, and each next-layer bit is a fixed Boolean function of the constant-size local predecessor window.  Introduce a universally quantified bit \(g_v\) for the output of every circuit gate.  Let \(\operatorname{Bad}_v\) be a disjunction of conjunctions, each of width at most three, listing the assignments on the input and output bits of gate \(v\) that violate its truth table.  For example, for \(g_v=g_a\wedge g_b\), the three bad terms are
\[
 g_a g_b\neg g_v,\qquad \neg g_a g_v,\qquad \neg g_b g_v.
\]
Negation, copy, constant, and input gates have analogous terms of width at most two.  Define
\[
 \Phi_w(x,z,g)=g_{\mathrm{out}}\ \vee\!\bigvee_v\operatorname{Bad}_v.
\]
This is a \(3\)-DNF of polynomial length.  For fixed \((x,z)\), every inconsistent assignment of \(g\) satisfies some bad term, while the unique consistent assignment satisfies \(\Phi_w\) exactly when \(C_w(x,z)=1\).  Consequently
\[
 w\in A
 \quad\Longleftrightarrow\quad
 \exists x\ \forall(z,g):\Phi_w(x,z,g)=1.
\]
Truth of one-alternation \(3\)-DNF formulas is therefore \(\Sigma_2^{\mathrm P}\)-hard; its membership follows by evaluating the matrix after the two guesses.  By the corrected \(\mathrm{lem{:}qbf\text{-}compiler\text{-}audit}\), \(\Phi\mapsto\mathcal C(\Phi)\) is a polynomial many-one reduction and
\[
 \exists x\,\forall z\,\Phi(x,z)=1
 \quad\Longleftrightarrow\quad
 \mathcal C(\Phi)\in\mathsf{RSJID}.
\]
The same lemma proves unit costs, singleton terminals, binary observed variables, two-mode signed guards, and maximum nonterminal bidirected degree eight.  It also shows explicitly that selector and universal coordinates can guard a nonconstant number of edge copies, so the proof establishes exactly the narrowed restriction package and no bounded-fanout claim.  This proves \(\Sigma_2^{\mathrm P}\)-hardness and hence completeness.

For the causal equivalence, fix a valid \(\mathbf i\), its family \(\mathfrak G\), and a budget-feasible \(J\).  The verifier lemma and \(\mathrm{thm{:}star\text{-}source\text{-}iff}\) give, for every released record \(h\),
\[
 \mathsf{Ver}(\mathbf i,J,h)=1
 \quad\Longleftrightarrow\quad
 T\text{ and }Y\text{ are disconnected in }H_h-J
 \quad\Longleftrightarrow\quad
 \theta_t(y)\text{ is identifiable from }\mathcal I_{h,J}
 \text{ over }\mathcal P_{+,h}(G_h,J).
\]
Here the middle notation means that no undirected bidirected-edge path survives.  Taking the universal quantifier over \(h\), and using the stipulated order in which \(J\) and \(\mathcal S_J\) precede release of \(h\), yields
\[
 \mathbf i\in\mathsf{RSJID}
 \quad\Longleftrightarrow\quad
 \mathcal J_B(\mathfrak G)\ne\varnothing.
\]
The completion-specific identifying map is the already derived \(\Psi_{h,J,t,y}\); no single unrevealed-graph formula is being substituted for it.
\end{proof}
\par\smallskip\noindent \textit{Justification.} It isolates a quantified complexity jump caused by committing to one experiment before correlated local graph modes are learned. \textit{Gap.} Akbari et al. and Elahi et al. study one known ADMG, while Bachtler et al. select forbidden arc pairs in a fixed graph; none proves this recorded-configuration causal vertex-intervention language under the restriction package. \textit{Consumer.} A dosearch planner can justify a QBF backend while retaining completion-specific gID semantics after release of the logged configuration record.

\begin{proposition}[prop:certifying-interface]\label{prop:certifying-interface}
For any supplied \(J\), released exogenous configuration record \(h\), and source \(\mathcal S_J\), the certificate interface runs in time polynomial in \(|V|+|E_h^{\leftrightarrow}|\). It returns \(D_h(J)\), \(A_h(J)\), and \(\Psi_{h,J,t,y}\) when identification holds over \(\mathcal P_{+,h}(G_h,J)\); otherwise it returns a simple path, its completion-specific \(Y\)-rooted hedge, and two same-fiber positive binary countermodels. Thus a robust \(J\) admits one polynomially checkable gID certificate per queried configuration, whereas a fixed failing \(J\) admits a configuration-plus-hedge witness.
\end{proposition}
\begin{proof}
Breadth-first search in \(H_h-J\) runs in \(O(|V|+|E_h^{\leftrightarrow}|)\) time.  If it does not reach \(T\) from \(Y\), its reached set is exactly \(D_h(J)\), so deleting \(Y\) gives \(A_h(J)\).  The expression defining \(\Psi_{h,J,t,y}\) can then be written from those vertex labels in polynomial time, and \(\mathrm{thm{:}star\text{-}source\text{-}iff}\) proves that it identifies the target over the fixed fiber.

If the search reaches \(T\), predecessor pointers return a simple path \(\pi=(T,A_1,\ldots,A_k,Y)\).  On the vertices of this path, take the outer forest with directed edges \(A_i\to T\) and \(T\to Y\), bidirected path edges from \(\pi\), and root \(Y\); every nonroot has one directed child leading to \(Y\), and the bidirected part is connected.  Its inner forest is the singleton \(\{Y\}\).  The outer forest contains \(T\), the inner one does not, and both survive after deleting the intervention vertices \(J\), so this pair is the completion-specific \(Y\)-rooted hedge carried by the path.  Independently of that graphical label, \(\mathrm{lem{:}path\text{-}countermodels}\) supplies the decisive semantic certificate: its two dyadic parameter tables define positive binary models in the same fiber with identical sources and different targets.  The integer \(k\), the rational \(2^{-k}\), and all table entries require only polynomially many bits.

Thus the success output is checked by reachability plus the explicit adjustment derivation, and the failure output is checked by the path, forest inclusions, and the displayed finite parameter calculation.  Applying the success branch separately to every queried completion gives a polynomially checkable certificate for each record.  For a fixed nonrobust \(J\), a completion \(h\) together with the returned path and hedge is a polynomial-size failure witness.  No claim is made that finding such an \(h\) without it being supplied is polynomial time.
\end{proof}
\par\smallskip\noindent \textit{Justification.} Both success and failure are checked relative to the fixed completion indexed by the same logged record that indexes the identifying functional. \textit{Gap.} Fixed-graph gID and PAG algorithms provide local proof objects, but the local-switch design literature lacks this fiber-correct two-sided interface. \textit{Consumer.} The objects feed directly into a dosearch pre-data design report and a QBF counterexample viewer.

\begin{proposition}[prop:fixed-completion-cut]\label{prop:fixed-completion-cut}
If \(R=\varnothing\), then \(\mathcal J_B(\mathfrak G)\ne\varnothing\) if and only if the minimum internal \(T\)-to-\(Y\) vertex-cut size in the fixed bidirected graph is at most \(B\); whenever this set is nonempty, its minimum-cardinality member is such a minimum cut. The cut and the feasibility decision are computable in polynomial time by vertex splitting and maximum flow, with value \(+\infty\) when an undeletable direct \(T\)-to-\(Y\) edge is present.
\end{proposition}
\begin{proof}
When \(R=\varnothing\), there is one completion and one residual bidirected graph \(H\).  By \(\mathrm{thm{:}star\text{-}source\text{-}iff}\), a set \(J\subseteq W\) belongs to \(\mathcal J_B(\mathfrak G)\) exactly when \(|J|\le B\) and \(J\) is an internal \(T\)-to-\(Y\) vertex separator in \(H\).  Hence feasibility is equivalent to the minimum separator size being at most \(B\), and, when feasible, minimizing \(|J|\) over \(\mathcal J_B(\mathfrak G)\) gives the unrestricted minimum separator.

For the algorithm, split every \(w\in W\) into \(w_{\mathrm{in}}\) and \(w_{\mathrm{out}}\), join them by a directed arc of capacity one, and replace every undirected edge by the two corresponding directed arcs of capacity \(M=|W|+1\).  Leave \(T\) and \(Y\) unsplit.  Any cut of capacity below \(M\) consists only of split arcs and corresponds, with equal capacity, to deletion of the associated proxy vertices.  Conversely, every proxy separator yields such a directed cut.  The cited minimum-weight vertex-cut result guarantees that this split-network computation returns the optimum in polynomial time.  Because all capacities are integral and any finite value is at most \(|W|\), an augmenting-path implementation makes at most \(|W|\) augmentations, each found in polynomial time.  If an undeletable direct edge joins \(T\) and \(Y\), every directed cut has capacity at least \(M\), which records the value \(+\infty\) for the allowed internal-separator problem.  This proves the claimed polynomial algorithm and the boundary case.
\end{proof}
\par\smallskip\noindent \textit{Justification.} This locates hardness in the ex ante universal completion quantifier rather than a hidden fixed-graph causal obstruction. \textit{Gap.} Known-graph intervention design is generally NP-complete, but the causal-star subclass collapses to ordinary minimum cut. \textit{Consumer.} The planner can dispatch single-completion stars to maximum flow.

\begin{proposition}[prop:greatest-completion-collapse]\label{prop:greatest-completion-collapse}
If there exists \(h^*\in\mathbb B^R\) such that \(E_h^{\leftrightarrow}\subseteq E_{h^*}^{\leftrightarrow}\) for every \(h\), then \(J\in\mathcal J_B(\mathfrak G)\) if and only if \(J\) disconnects \(T\) and \(Y\) in \(H_{h^*}-J\); equivalently, identification in the fixed fiber at \(h^*\) certifies identification in every fixed fiber.
\end{proposition}
\begin{proof}
Assume the greatest completion \(h^*\) exists.  Vertex deletion preserves edge inclusion, so for every \(J\subseteq W\) and every \(h\),
\[
 E(H_h-J)\subseteq E(H_{h^*}-J).
\]
If \(J\) disconnects \(T\) and \(Y\) in \(H_{h^*}-J\), then a path in any \(H_h-J\) would also be a path in \(H_{h^*}-J\), a contradiction.  Thus it disconnects every completion.  Conversely, separation in every completion includes separation at \(h^*\).  Applying \(\mathrm{thm{:}star\text{-}source\text{-}iff}\) separately at each fixed record and retaining the same budget condition yields
\[
 J\in\mathcal J_B(\mathfrak G)
 \quad\Longleftrightarrow\quad
 |J|\le B\text{ and }T,Y\text{ are disconnected in }H_{h^*}-J.
\]
The right side is, again by that theorem, identification in the fixed fiber at \(h^*\).  This proves the collapse without taking a union of source laws across different fibers.
\end{proof}
\par\smallskip\noindent \textit{Justification.} Complementary correlations, rather than uncertainty with a largest completion, cause the quantified hardness. \textit{Gap.} Robust-flow models often admit a worst scenario; this proposition identifies the exact collapse for the present completion-indexed source protocol. \textit{Consumer.} A planner can replace robust search by one minimum cut whenever a greatest completion exists.

\begin{theorem}[thm:bounded-terminal-frontier-complete]\label{thm:bounded-terminal-frontier-complete}
In the single-signed-coordinate guard model defining \(\mathsf{RSJID}\), if every completion has bidirected degree at most a fixed \(d\) at \(T\) and at \(Y\), then the union of all fixed and possible bidirected edges has degree at most \(2d\) at each terminal. Thus the proposed regime with fixed completion-wise terminal degree but unbounded union terminal neighborhoods is empty. On the corresponding promise restriction of \(\mathsf{RSJID}\), the decision problem is trivial for \(d=0\), and it is \(\mathsf{coNP}\)-complete for every fixed \(d\ge1\), already for \(B=0\) and completion-wise terminal degree one. This classification does not additionally impose a bounded guard-fanout promise.
\end{theorem}
\begin{proof}
Fix a terminal, say \(T\).  Let \(f\) be the number of fixed bidirected edges incident to \(T\), and for each \(r\in R\) let \(n_r^+\) and \(n_r^-\) be the numbers of possible incident edges whose guards are \((r,+1)\) and \((r,-1)\), respectively.  Because the coordinates of \(h\) are independently assignable and each possible edge has one signed coordinate, the maximum completion degree at \(T\) is
\[
 f+\sum_{r\in R}\max\{n_r^+,n_r^-\}.
\]
The assumed completion-wise bound therefore implies
\[
 \deg_{E^{\mathrm{fix}}\cup E^{\cup}}(T)
 =f+\sum_{r\in R}(n_r^++n_r^-)
 \le f+2\sum_{r\in R}\max\{n_r^+,n_r^-\}
 \le 2d-f\le2d.
\]
The same argument applies to \(Y\).  In particular, an unbounded union neighborhood is impossible for fixed \(d\) under this guard semantics.

We next prove the upper bound.  Let \(S\) be the union of all possible bidirected neighbors of \(T\), including fixed neighbors.  The preceding inequality gives \(|S|\le2d\).  If \(Y\in S\), some completion contains the undeletable direct edge \(T\leftrightarrow Y\), and the instance is a no-instance.  Otherwise \(S\subseteq W\), and deleting all of \(S\) isolates \(T\) in every completion.  Thus the instance is immediately a yes-instance if \(B\ge|S|\).  In the remaining case \(B<|S|\le2d\).  For fixed \(d\), there are only
\[
 \sum_{b=0}^{B}\binom{|W|}{b}=O(|W|^{2d})
\]
candidate sets.  The complement of robust feasibility has an NP certificate: for every one of these polynomially many candidates \(J\), list a completion \(h_J\) for which breadth-first search finds a residual \(T\)-to-\(Y\) path.  Each record has length \(|R|\), and \(\mathrm{lem{:}fixed\text{-}completion\text{-}verifier}\) checks every listed witness in polynomial time.  Therefore the fixed-\(d\) promise restriction lies in \(\mathsf{coNP}\).  When \(d=0\), \(T\) has no incident bidirected edge in any completion, so \(J=\varnothing\) always works and the valid-input promise problem is trivial.

For hardness at every \(d\ge1\), reduce from the cited \(3\)-CNF unsatisfiability language.  Given a \(3\)-CNF formula with clauses \(C_1,\ldots,C_m\), create clause-boundary vertices \(q_0,\ldots,q_m\), a fixed edge \(T\leftrightarrow q_0\), and a fixed edge \(q_m\leftrightarrow Y\).  For every literal occurrence \(\ell\in C_i\), create a private vertex \(a_{i,\ell}\), the fixed edge \(q_{i-1}\leftrightarrow a_{i,\ell}\), and the edge \(a_{i,\ell}\leftrightarrow q_i\) guarded by the corresponding variable coordinate with positive sign for a positive literal and negative sign for a negated literal.  Set \(B=0\).  Equality and complement copies are exactly repeated positive and negative uses of one coordinate.  In a completion \(h\), the \(i\)-th clause block connects \(q_{i-1}\) to \(q_i\) exactly when \(C_i(h)\) is true.  The blocks are in series, so
\[
 T\text{ is connected to }Y
 \quad\Longleftrightarrow\quad
 \bigwedge_{i=1}^m C_i(h)=1.
\]
Because \(B=0\), the only design is empty.  It is robust exactly when no completion satisfies the formula, that is, exactly when the input formula is unsatisfiable.  Both terminals have the single fixed bidirected neighbor \(q_0\) or \(q_m\) in every completion, so their completion-wise degrees are one.  The construction is polynomial and uses the required single-coordinate equality/complement guards.  This proves \(\mathsf{coNP}\)-hardness for every fixed \(d\ge1\), and the upper bound proves completeness.
\end{proof}

\begin{lemma}[lem:sigma-two-verifier-characterization]\label{lem:sigma-two-verifier-characterization}
A language lies in \(\Sigma _2^{\mathrm P}\) if and only if membership has an \(\exists x\,\forall z\) characterization by a deterministic polynomial-time predicate, with both quantified strings of polynomial length.
\end{lemma}

\begin{lemma}[lem:max-flow-min-cut]\label{lem:max-flow-min-cut}
\texttt{The minimum-\allowbreak{}weight internal vertex-\allowbreak{}cut problem on a finite directed or undirected graph with nonnegative vertex weights is solvable in polynomial time by vertex splitting followed by a maximum-\allowbreak{}flow computation.}
\end{lemma}

\begin{lemma}[lem:three-cnf-unsat-complete]\label{lem:three-cnf-unsat-complete}
The language of unsatisfiable Boolean formulas in \(3\)-CNF is \(\mathsf{coNP}\)-complete under polynomial many-one reductions.
\end{lemma}

\section*{Technical internal limitation (diagnostic only)}

The compiler audit remains a proof obligation: bounded checks support, but do not replace, the uniform arguments for occurrence bookkeeping, selector isolation, and the rounding map. Exact dyadic checks for paths with zero through four internal proxies support, but do not replace, the symbolic same-fiber induction. Terminal degrees remain unbounded.

\section{Honest scope}

Identification is always asserted within a fixed fiber \(\mathcal P_{+,h}(G_h,J)\) after the directly logged exogenous configuration record \(h\) is released with \(\mathcal S_J\); latent bidirected edges are not inferred from the acquired laws. No identification claim is made over an unindexed union of completions or for a formula chosen before metadata release. The result is restricted to directed causal stars and strictly positive binary models. The complexity theorem bounds nonterminal bidirected degree but not terminal degree. No claim is made for PAG families, unrestricted strictly binary gID completeness, finite-sample estimation, or polynomial-time discovery of a failing completion. With \(R=\varnothing\) the design problem is ordinary minimum vertex cut; a greatest completion suffices when present; and if \(A_h(J)=\varnothing\), the functional reduces to \(Q_{J,j_0}(Y=y\mid T=t)\).

\begin{thebibliography}{99}

  \bibitem{akbari2023uncertain} Sina Akbari, Fateme Jamshidi, Ehsan Mokhtarian, Matthew J. Vowels, Jalal Etesami, and Negar Kiyavash. Causal Effect Identification in Uncertain Causal Networks. Advances in Neural Information Processing Systems 36, 2023. arXiv:2208.04627.

  \bibitem{akbari2025optimal} Sina Akbari, Jalal Etesami, and Negar Kiyavash. Optimal Experiment Design for Causal Effect Identification. Journal of Machine Learning Research 26(28):1--56, 2025. arXiv:2205.02232.

  \bibitem{elahi2024fast} Sepehr Elahi, Sina Akbari, Jalal Etesami, Negar Kiyavash, and Patrick Thiran. Fast Proxy Experiment Design for Causal Effect Identification. Advances in Neural Information Processing Systems 37, 2024. arXiv:2407.05330.

  \bibitem{kivva2022revisiting} Yaroslav Kivva, Ehsan Mokhtarian, Jalal Etesami, and Negar Kiyavash. Revisiting the General Identifiability Problem. Proceedings of Machine Learning Research 180:1022--1030, UAI, 2022. arXiv:2206.01081.

  \bibitem{lee2020general} Sanghack Lee, Juan D. Correa, and Elias Bareinboim. General Identifiability with Arbitrary Surrogate Experiments. Proceedings of Machine Learning Research 115:389--398, UAI, 2020.

  \bibitem{jaber2022causal} Amin Jaber, Adèle H. Ribeiro, Jiji Zhang, and Elias Bareinboim. Causal Identification under Markov Equivalence: Calculus, Algorithm, and Completeness. Advances in Neural Information Processing Systems 35:3679--3690, 2022.

  \bibitem{bachtler2024almost} Oliver Bachtler, Tim Bergner, and Sven O. Krumke. Almost Disjoint Paths and Separating by Forbidden Pairs. Theoretical Computer Science 982:114272, 2024.

  \bibitem{gruenewulf2026blocking} Christoph Grüne and Lasse Wulf. The Complexity of Blocking All Solutions. Theoretical Computer Science 1069:115820, 2026.

  \bibitem{bertsimas2013robust} Dimitris Bertsimas, Ebrahim Nasrabadi, and Sebastian Stiller. Robust and Adaptive Network Flows. Operations Research 61(5):1218--1242, 2013.

  \bibitem{rebennack2020twostage} Steffen Rebennack, Oleg A. Prokopyev, and Bismark Singh. Two-Stage Stochastic Minimum Cut Problems: Formulations, Complexity and Decomposition Algorithms. Networks 75(3):235--258, 2020.

  \bibitem{shpitserpearl2006} Ilya Shpitser and Judea Pearl. Identification of Joint Interventional Distributions in Recursive Semi-Markovian Causal Models. Proceedings of AAAI, pages 1219--1226, 2006.

\end{thebibliography}

\end{document}